\documentclass{beamer}
\usetheme{Madrid}

\title{Impact of COVID-19 on the Labor Market in Canada}
\subtitle{Contemporary Economic Issues: ECON 204}
\author{Muhebullah Karimzada}
\date{Winter 2025}
\begin{document}


\frame{\titlepage}

\section{Introduction}
\begin{frame}{Introduction}
    \begin{itemize}
        \item COVID-19 profoundly disrupted Canada's labor market.
        \item Key points of the lecture:
        \begin{itemize}
            \item Immediate impacts
            \item Sectoral effects
            \item Demographic disparities
            \item Government responses
            \item Long-term consequences
            \item Policy implications
        \end{itemize}
    \end{itemize}
\end{frame}

\section{Immediate Impacts}
\begin{frame}{Immediate Impacts: Employment and Hours Worked}
    \begin{itemize}
        \item \textbf{Employment losses:} 
        \begin{itemize}
            \item 3 million jobs lost (March–April 2020)
            \item Unemployment peaked at 13.7\% in May 2020
        \end{itemize}
        \item \textbf{Hours worked:}
        \begin{itemize}
            \item 28\% decline in total hours worked
        \end{itemize}
    \end{itemize}
    \textbf{Case Study: Ontario's Manufacturing Sector}
    \begin{itemize}
        \item 45\% drop in manufacturing employment during lockdowns.
        \item Ford retooled plants to manufacture ventilators.
    \end{itemize}
\end{frame}

\begin{frame}{Nature of Job Losses}
    \begin{itemize}
        \item Nearly 50\% of job losses were in temporary or part-time positions.
        \item Non-essential services hit hardest (retail, hospitality, arts).
    \end{itemize}
\end{frame}

\section{Sectoral Effects}
\begin{frame}{Hardest-hit Sectors}
    \textbf{Hospitality and Tourism}
    \begin{itemize}
        \item Hotel occupancy rates fell below 10\% (April 2020).
        \item Example: Fairmont Hotels layed off over 3,000 workers.
    \end{itemize}
    \textbf{Arts and Entertainment}
    \begin{itemize}
        \item Example: Cirque du Soleil laid off 95\% of its workforce.
    \end{itemize}
    \textbf{Retail Sector}
    \begin{itemize}
        \item Small businesses in Toronto's downtown core faced closures.
        \item E-commerce surged: Shopify reported a 76\% increase in new online stores.
    \end{itemize}
\end{frame}

\begin{frame}{Resilient or Growing Sectors}
    \begin{itemize}
        \item \textbf{Technology and e-commerce:}
        \begin{itemize}
            \item Amazon hired 8,000 workers in 2020.
        \end{itemize}
        \item \textbf{Health care:}
        \begin{itemize}
            \item Example: British Columbia hired contact tracers, increasing capacity by 200\%.
        \end{itemize}
        \item \textbf{Essential services:}
        \begin{itemize}
            \item Supermarkets, pharmacies, and delivery services expanded their workforce.
        \end{itemize}
    \end{itemize}
\end{frame}

\section{Demographic Disparities}
\begin{frame}{Gender and Youth Disparities}
    \textbf{Gender Disparities}
    \begin{itemize}
        \item Women accounted for 63\% of job losses.
        \item Many women left the workforce due to caregiving responsibilities.
    \end{itemize}
    \textbf{Youth Employment}
    \begin{itemize}
        \item Unemployment for ages 15–24 peaked at 29.4\%.
        \item Example: Canada Summer Jobs Program expanded to subsidize 70,000 placements.
    \end{itemize}
\end{frame}

\begin{frame}{Income Inequality}

    \textbf{Income Inequality}
    \begin{itemize}
        \item Low-wage workers faced the steepest losses.
        \item High-income workers in remote-capable jobs recovered quickly.
    \end{itemize}
\end{frame}

\section{Government Response}
\begin{frame}{Emergency Benefits}
    \begin{itemize}
        \item \textbf{CERB:} $\$2,000/month$ to 8 million Canadians.
        \item \textbf{CEWS:} Subsidized wages for 4 million jobs.
    \end{itemize}
    \textbf{Sector-specific Aid}
    \begin{itemize}
        \item Airlines and tourism received $\$5$ billion in loans.
    \end{itemize}
\end{frame}

\section{Long-term Consequences}
\begin{frame}{Structural Changes and Mental Health}
    \begin{itemize}
        \item \textbf{Remote Work:} Hybrid models became the norm for many sectors.
        \item \textbf{Automation:} Manufacturing firms accelerated automation.
        \item \textbf{Mental Health:} Burnout among healthcare workers rose sharply.
    \end{itemize}
\end{frame}

\section{Policy Implications}
\begin{frame}{Policy Implications}
    \begin{itemize}
        \item Strengthen social safety nets: Flexible, inclusive programs.
        \item Invest in skills training: Support for transitioning into high-demand fields.
        \item Address inequality: Subsidized childcare and equity-focused programs.
        \item Prepare for future shocks: Paid sick leave, workplace digitization.
        \item Focus on mental health: Workplace benefits and systemic reforms.
    \end{itemize}
\end{frame}

\section{Conclusion}
\begin{frame}{Conclusion}
    \begin{itemize}
        \item The pandemic revealed vulnerabilities and opportunities in Canada's labour market.
        \item Collaborative policies can build a more resilient and equitable economy.
    \end{itemize}
 
\end{frame}




% Problem 1
\begin{frame}{Problem 1: Dynamic Unemployment Rate Analysis}
In May 2020, the labor force in Canada was 19 million people, and 2.6 million were unemployed. By September 2020, 1 million unemployed workers returned to work, but 500,000 people dropped out of the labor force.

\begin{enumerate}
    \item Calculate the unemployment rate in May 2020.
    \item Calculate the unemployment rate in September 2020, accounting for the reduced labor force.
    \item By how much did the unemployment rate decrease between May and September?
\end{enumerate}
\end{frame}

\begin{frame}{Solution: Problem 1}
\begin{enumerate}
    \item Unemployment rate (May 2020): 
    \[
    \text{Unemployment Rate} = \frac{\text{Unemployed}}{\text{Labor Force}} \times 100 = \frac{2.6M}{19M} \times 100 = 13.68\%.
    \]
    \item Reduced labor force: 
    \[
    \text{New Labor Force} = 19M - 0.5M = 18.5M.
    \]
    Unemployed workers: 
    \[
    \text{New Unemployed} = 2.6M - 1M = 1.6M.
    \]
    Unemployment rate (September 2020): 
    \[
    \frac{1.6M}{18.5M} \times 100 = 8.65\%.
    \]
    \item Decrease in unemployment rate: 
    \[
    13.68\% - 8.65\% = 5.03\% \text{ (percentage points)}.
    \]
\end{enumerate}
\end{frame}

% Problem 2
\begin{frame}{Problem 2: Sectoral Employment Decline and Recovery}
The hospitality sector employed 1.5 million workers in January 2020. During the first wave, employment dropped by 45\%. By December 2021, employment recovered by 60\% of the initial job losses.

\begin{enumerate}
    \item How many workers were employed at the peak of the job losses?
    \item How many jobs were recovered by December 2021?
    \item What percentage of the pre-pandemic workforce was employed by December 2021?
\end{enumerate}
\end{frame}

\begin{frame}{Solution: Problem 2}
\begin{enumerate}
    \item Jobs lost: 
    \[
    0.45 \times 1.5M = 675,000.
    \]
    Workers remaining: 
    \[
    1.5M - 675,000 = 825,000.
    \]
    \item Jobs recovered: 
    \[
    0.60 \times 675,000 = 405,000.
    \]
    Workers employed: 
    \[
    825,000 + 405,000 = 1,230,000.
    \]
    \item Percentage of pre-pandemic workforce: 
    \[
    \frac{1,230,000}{1,500,000} \times 100 = 82\%.
    \]
\end{enumerate}
\end{frame}

% Problem 3
\begin{frame}{Problem 3: CERB and Household Income Analysis}
The Canadian government provided CERB payments of \$2,000/month to 8 million people for 6 months. Suppose the average pre-pandemic household income for these recipients was \$3,500/month.

\begin{enumerate}
    \item Calculate the total CERB expenditure over 6 months.
    \item What percentage of their pre-pandemic income did CERB replace?
    \item If household expenses averaged \$2,700/month, calculate the monthly shortfall or surplus during CERB.
\end{enumerate}
\end{frame}

\begin{frame}{Solution: Problem 3}
\begin{enumerate}
    \item Total CERB expenditure: 
    \[
    8M \times \$2,000 \times 6 = \$96 \text{ billion}.
    \]
    \item Percentage of income replaced: 
    \[
    \frac{\$2,000}{\$3,500} \times 100 = 57.14\%.
    \]
    \item Monthly shortfall: 
    \[
    \$2,700 - \$2,000 = \$700 \text{ (shortfall)}.
    \]
\end{enumerate}
\end{frame}

% Problem 4
\begin{frame}{Problem 4: Wage Subsidy Program and Business Payrolls}
A company applied for CEWS to cover 75\% of payroll costs for 10 employees earning \$4,000/month each. The program lasted for 6 months.

\begin{enumerate}
    \item Calculate the total subsidy the company received.
    \item If the company’s revenue was \$300,000 during the 6 months, what percentage of its revenue was used to cover the remaining payroll costs?
    \item Assuming the business’s total operating costs (excluding payroll) were \$120,000, calculate the profit or loss during the 6-month period.
\end{enumerate}
\end{frame}

\begin{frame}{Solution: Problem 4}
\begin{enumerate}
    \item Total subsidy: 
    \[
    75\% \times (10 \times \$4,000) \times 6 = \$180,000.
    \]
    \item Remaining payroll: 
    \[
    25\% \times (10 \times \$4,000) \times 6 = \$60,000.
    \]
    Percentage of revenue: 
    \[
    \frac{\$60,000}{\$300,000} \times 100 = 20\%.
    \]
    \item Profit or loss: 
    \[
    \text{Revenue - Total Costs} = \$300,000 - (\$60,000 + \$120,000) = \$120,000 \text{ profit}.
    \]
\end{enumerate}
\end{frame}



% Problem 5
\begin{frame}{Problem 5: E-commerce Sector Growth}
The e-commerce sector grew by 76\% in 2020 compared to pre-pandemic levels. In 2019, the sector generated \$50 billion in revenue. By 2021, growth slowed to 12\% compared to 2020.

\begin{enumerate}
    \item Calculate the revenue in 2020.
    \item Calculate the revenue in 2021.
    \item What was the compound annual growth rate (CAGR) from 2019 to 2021?
\end{enumerate}
\end{frame}

\begin{frame}{Solution: Problem 5}
\begin{enumerate}
    \item Revenue in 2020:
    \[
    50B \times (1 + 0.76) = 50B \times 1.76 = 88B.
    \]
    \item Revenue in 2021:
    \[
    88B \times (1 + 0.12) = 88B \times 1.12 = 98.56B.
    \]
    \item CAGR:
    \[
    \text{CAGR} = \left(\frac{\text{Revenue in 2021}}{\text{Revenue in 2019}}\right)^{\frac{1}{2}} - 1
    \]
    Substituting values:
    \[
    \text{CAGR} = \left(\frac{98.56}{50}\right)^{\frac{1}{2}} - 1 = (1.9712)^{0.5} - 1 = 1.396 - 1 = 0.396 = 39.6\%.
    \]
\end{enumerate}
\end{frame}

% Problem 6
\begin{frame}{Problem 6: Regional Tourism Losses}
In 2020, tourism in British Columbia generated \$10 billion less than the previous year, representing a 60\% decline. Assume tourism gradually recovers, growing 20\% annually starting in 2021.

\begin{enumerate}
    \item What was the total tourism revenue in 2019?
    \item Calculate the tourism revenue in 2021 and 2022.
    \item How many years will it take to fully recover to 2019 levels?
\end{enumerate}
\end{frame}

\begin{frame}{Solution: Problem 6}
\begin{enumerate}
    \item Revenue in 2019:
    \[
    \text{Revenue in 2019} = \frac{\$10B}{0.6} = 16.67B.
    \]
    \item Revenue in 2021:
    \[
    10B \times (1 + 0.20) = 10B \times 1.2 = 12B.
    \]
    Revenue in 2022:
    \[
    12B \times (1 + 0.20) = 12B \times 1.2 = 14.4B.
    \]
    \item Time to recover:
    \[
    10B \times (1.20)^x = 16.67B.
    \]
    Taking the natural logarithm:
    \[
    x = \frac{\ln(1.667)}{\ln(1.20)} \approx 2.8.
    \]
    Full recovery will take approximately 3 years.
\end{enumerate}
\end{frame}

% Problem 7
\begin{frame}{Problem 7: Labor Force Dynamics}
In 2020, the working-age population in Canada was 30 million. The labor force participation rate was 65\%, and the employment-to-population ratio was 58\%.

\begin{enumerate}
    \item Calculate the number of people in the labor force and the number employed.
    \item If 500,000 workers dropped out of the labor force in 2021, calculate the new participation rate (assume the working-age population remains constant).
    \item If employment increased by 800,000 in 2021, calculate the new employment-to-population ratio.
\end{enumerate}
\end{frame}

\begin{frame}{Solution: Problem 7}
\begin{enumerate}
    \item Labor force:
    \[
    \text{Labor Force} = 0.65 \times 30M = 19.5M.
    \]
    Number employed:
    \[
    \text{Employed} = 0.58 \times 30M = 17.4M.
    \]
    \item New participation rate:
    \[
    \text{New Labor Force} = 19.5M - 0.5M = 19M.
    \]
    \[
    \text{Participation Rate} = \frac{19M}{30M} \times 100 = 63.33\%.
    \]
    \item New employment-to-population ratio:
    \[
    \text{New Employment} = 17.4M + 0.8M = 18.2M.
    \]
    \[
    \text{Employment-to-Population Ratio} = \frac{18.2M}{30M} \times 100 = 60.67\%.
    \]
\end{enumerate}
\end{frame}

\end{document}
